% SynAthlete-AIoT project report (Review 2 rubric: "Project report & draft of paper").
% Same rule as the paper: every number comes from generated/*.tex (python tasks.py paper-figures).
\documentclass[11pt,a4paper]{report}
\usepackage[margin=2.5cm]{geometry}
\usepackage{graphicx,booktabs,amsmath,url,xspace,hyperref}
% generated by src/paper/figures.py: do not edit
\newcommand{\canaryAuc}{0.498 (day) and 0.513 (week)\xspace}
\newcommand{\explainSummary}{Explanation checks: n/a.\xspace}
\newcommand{\fusedAuc}{n/a\xspace}
\newcommand{\fusedPrAuc}{n/a\xspace}
\newcommand{\injuryAucRange}{0.58--0.63\xspace}
\newcommand{\labelFreeAuc}{0.584 (day) and 0.589 (week)\xspace}
\newcommand{\lockboxSummary}{Pending the evaluation freeze.\xspace}
\newcommand{\lovdalRunners}{74\xspace}
\newcommand{\pamapRealFone}{0.632\xspace}
\newcommand{\reproDayAuc}{0.617 (95\% CI 0.573--0.664)\xspace}
\newcommand{\tiersReduction}{n/a\xspace}

\title{SynAthlete-AIoT\\[4pt]\large An Explainable AIoT Framework for Athlete Injury-Risk
Prediction Using Synthesized Wearable Data\\[10pt]\normalsize Project Report (Review 2),
21CSE292P Artificial Intelligence of Things (PBL)}
\author{[Team names]\\ SRM Institute of Science and Technology}
\date{October 2026}

\begin{document}
\maketitle
\tableofcontents

\chapter{Introduction}
\section{Problem}
Coaches cannot see injury risk building up over days. Wearable AIoT systems promise early
warning, but published athlete-health frameworks mostly describe the architecture and leave
the prediction and explanation layers unimplemented~\cite{aiotframework2026}.
\section{Research gap (corrected since Review 1)}
L\"ovdal et al.\ published seven years of daily training logs from \lovdalRunners{} competitive
runners with real injuries~\cite{lovdal2021data}, so the gap is narrower than first stated:
\emph{no public dataset pairs high-frequency physiological streams with real injury outcomes}.
\section{Objectives}
\begin{enumerate}
\item Bridge the data gap with synthetic wearable data (TimeGAN, WGAN-GP; TimeVAE and classic
augmentation as comparisons).
\item Implement simulated edge--fog--cloud tiers with measured bandwidth and latency.
\item Build a fused model (sequence + tabular arms) producing a calibrated 0--100 risk score.
\item Explain predictions with SHAP and attribution methods, validated by tests.
\item Deliver a coach dashboard.
\end{enumerate}
\section{Honesty rules}
Models are tested only on people they never saw; injury claims come only from real injury
labels; synthetic data is labelled synthetic and used only for training; every number in this
report is produced by code.

\chapter{Literature Review}
Injury prediction from training load has been studied with tree ensembles and boosting;
L\"ovdal et al.\ used bagged XGBoost~\cite{lovdal2021,chen2016}. Tree models remain strong on
tabular data~\cite{grinsztajn2022,breiman2001,prokhorenkova2018}, recurrent networks read
sequences~\cite{cho2014}, and convolutional and recurrent models recognise activities from
wearables~\cite{hammerla2016}. TimeGAN~\cite{yoon2019}, WGAN-GP~\cite{gulrajani2017} and
TimeVAE~\cite{desai2021} generate time series. SHAP~\cite{lundberg2017} and integrated
gradients~\cite{sundararajan2017} explain predictions, but must pass sanity
checks~\cite{adebayo2018}. Leakage and selection outside the test loop inflate machine-learning
results~\cite{kapoor2023,varma2006,cawley2010}.

\chapter{Datasets and Their Roles}
Every dataset was audited (origin, licence, contents) and given one role
(Table~\ref{tab:data}). Only the L\"ovdal logs carry real injuries; the wearable datasets
support the IoT and synthetic-data parts; datasets of unknown origin are demos only.
\begin{table}[t]\centering\caption{Datasets and their single role}\label{tab:data}
\footnotesize\begin{tabular}{p{2.6cm}p{1.2cm}p{3.6cm}}\toprule
Dataset & Origin & Role \\\midrule
Lövdal runners (74 runners, 7 years of daily logs) & real & the only injury evidence: baselines, combined model, explanations, coach view \\
Kaggle 'Injury Prediction for Competitive Runners' & real (copy) & byte-identical copy of Lövdal: no separate role \\
PAMAP2 (9 people, wearable sensors) & real & edge sensor stream, activity recognition, generator training and fidelity \\
WESAD (15 people, wrist sensors) & real & second dataset for generator generalisation \\
Sport DB 2.0 (55 heart-rate recordings) & real & heart-rate realism check and second edge stream \\
SoccerMon, PMData & real & exploratory only (Track C), never main evidence \\
SIRP-600 & unknown origin & demo only: shows the SHAP code runs \\
ziya07 wearable & unknown origin & demo only: dashboard demo stream \\
\bottomrule\end{tabular}\end{table}


\chapter{Methodology}
\section{Evaluation protocol}
Splits are by person and were created once. The ten most recently joined runners are a sealed
final test, opened once after an evaluation freeze. The remaining 64 runners form five
development rounds, each with training, validation and test runners; validation runners choose
settings, calibration, the alert level and early stopping; test runners are scored once.
PR-AUC is the headline score (rare events~\cite{saito2015}), reported next to its chance level;
ROC-AUC is reported for comparison. Uncertainty uses a runner-level cluster
bootstrap~\cite{field2007}.
\section{Leakage controls}
One scaler fitted on training runners only; inputs never include identities, dates or a
runner's own event history (which reveals the label in this dataset); a shuffled-label leak
alarm~\cite{ojala2010}; and an independent code review. All rules were written in the project's
decision log before the runs they govern.
\section{System design}
Data $\rightarrow$ synthetic data $\rightarrow$ edge (reflex) $\rightarrow$ fog (supervisor)
$\rightarrow$ cloud (deep model) $\rightarrow$ explanations $\rightarrow$ coach dashboard.

\chapter{Implementation}
\section{Code structure}
\begin{tabular}{ll}\toprule
Part & Code \\\midrule
Data loading and person-level splits & \texttt{src/data/}, \texttt{src/models/data.py} \\
Scoring rules and uncertainty & \texttt{src/eval/metrics.py}, \texttt{bootstrap.py} \\
Baselines and paper rebuild & \texttt{src/models/run.py}, \texttt{lovdal\_method.py} \\
Generators and their checks & \texttt{src/gen/} \\
Edge--fog--cloud tiers & \texttt{src/pipeline/} \\
Fused model & \texttt{src/models/fused.py} \\
Explanations and checks & \texttt{src/xai/} \\
Dashboard & \texttt{src/dashboard/} \\\bottomrule
\end{tabular}
\section{Baselines}
Constant risk, an ACWR-style rule~\cite{lolli2019}, logistic regression, random forest,
CatBoost, a GRU, and the L\"ovdal bagged XGBoost (exact copy for a reproduction check; label-free
variant for comparison). Eight settings per model and round by seeded random
search~\cite{bergstra2012}; the chosen setting is refit with three seeds.
\section{Synthetic data}
Class-conditional TimeGAN, WGAN-GP and TimeVAE are trained only on each round's training people;
settings are chosen by validation train-on-synthetic-test-on-real macro-F1; classic augmentation
is the non-learned comparison. Pass marks (realism, usefulness, privacy) were fixed in advance.
\section{Edge--fog--cloud}
Three operating-system processes connected by queues: the edge checks each 5-second window
(heart-rate and motion rule plus an activity CNN) and sends a short summary; the fog groups
windows into sessions; the cloud runs the daily injury model. A cloud-only mode sends raw windows.
\section{Fused model}
Sequence arm (GRU, LSTM or TCN~\cite{bai2018}) and tabular arm (random forest or CatBoost),
combined by logistic regression on out-of-fold predictions grouped by runner; the calibrated
output is shown as a 0--100 risk score.
\section{Explanations}
TreeSHAP, integrated gradients, an explanation-only attention model and 1D Grad-CAM, each
checked by deletion, model-randomisation and stability tests.
\section{Dashboard}
Eight Streamlit pages (home, athletes, athlete detail, live sensor, synthetic data, results,
training lab, data and honesty), built from data prepared in advance.

\chapter{Results}
\begin{table}[t]\centering\caption{Injury prediction, mean $\pm$ SD over five development rounds (chance PR-AUC $\approx$ injury rate)}\label{tab:injury}
\footnotesize\begin{tabular}{llcc}\toprule
Model & File & PR-AUC & ROC-AUC \\\midrule
Constant risk & day & 0.013 $\pm$ 0.000 & 0.500 $\pm$ 0.000 \\
Constant risk & week & 0.013 $\pm$ 0.000 & 0.500 $\pm$ 0.000 \\
ACWR-style rule & week & 0.015 $\pm$ 0.001 & 0.540 $\pm$ 0.021 \\
Logistic regression & day & 0.022 $\pm$ 0.004 & 0.626 $\pm$ 0.019 \\
Logistic regression & week & 0.020 $\pm$ 0.004 & 0.621 $\pm$ 0.029 \\
Random forest & day & 0.019 $\pm$ 0.002 & 0.603 $\pm$ 0.024 \\
Random forest & week & 0.022 $\pm$ 0.007 & 0.614 $\pm$ 0.042 \\
CatBoost & day & 0.021 $\pm$ 0.004 & 0.598 $\pm$ 0.043 \\
CatBoost & week & 0.019 $\pm$ 0.002 & 0.606 $\pm$ 0.028 \\
GRU & day & 0.023 $\pm$ 0.006 & 0.631 $\pm$ 0.009 \\
L\"ovdal, label-free & day & 0.017 $\pm$ 0.003 & 0.584 $\pm$ 0.044 \\
L\"ovdal, label-free & week & 0.021 $\pm$ 0.006 & 0.589 $\pm$ 0.066 \\
\textbf{Fused} & shared & 0.022 $\pm$ 0.003 & 0.609 $\pm$ 0.024 \\
\bottomrule\end{tabular}\end{table}

The leak alarm passed (shuffled-label ROC-AUC \canaryAuc{}). Learned injury models reach ROC-AUC
\injuryAucRange{}. The exact copy of the published method passed its reproduction check
(ROC-AUC \reproDayAuc{}); without its label-reading normalisation it reaches \labelFreeAuc{}.
The fused model reaches ROC-AUC \fusedAuc{}. Activity recognition on unseen people reaches
macro-F1 \pamapRealFone{}.
\begin{table}[t]\centering\caption{Generators on PAMAP2 (4 rounds): discriminative accuracy (0.5 = indistinguishable), TSTR and real+synthetic macro-F1, pre-registered criteria passed}\label{tab:gen}
\footnotesize\begin{tabular}{lcccc}\toprule
Generator & Disc.\ acc. & TSTR & Real+syn. & Passed \\\midrule
\multicolumn{5}{c}{n/a} \\
\bottomrule\end{tabular}\end{table}

\begin{table}[t]\centering\caption{Edge--fog--cloud vs cloud-only (simulated on one computer)}\label{tab:tiers}
\footnotesize\begin{tabular}{llrr}\toprule
Stream & Mode & Bytes into cloud & Median decision (ms) \\\midrule
\multicolumn{4}{c}{n/a} \\
\bottomrule\end{tabular}\end{table}

\explainSummary{} \lockboxSummary{}

\chapter{Testing and Validation}
The project has automated tests for every module (run with \texttt{python tasks.py check}):
split integrity and lockbox protection, scoring rules with hand-worked examples, leak guards
(the scaler sees training rows only; the combiner never sees predictions from models trained on
the same runner), generator shapes and seeding, tier message flow, explanation checks and every
dashboard page. Independent reviews of the evaluation code found no blocking problem; their
smaller findings were fixed and logged.

\chapter{Conclusion and Future Work}
The pipeline is implemented end to end and evaluated honestly. Injury risk from training logs is
only modestly predictable; the remaining work is the final locked test, the camera-ready paper
and the conference submission.

\appendix
\chapter{AI-Use Statement}
AI coding assistants (Claude Code; Antigravity with Claude Opus 4.6 and Gemini 3.8 Flash) helped
write code and text under human direction. Every AI contribution is logged in
\texttt{docs/AI\_LOG.md}; all numbers come from code; all claims were checked by the team.

\bibliographystyle{IEEEtran}
\bibliography{refs}
\end{document}
